\documentclass{article}
\usepackage{amssymb}
\usepackage{amsmath}
\usepackage{geometry}
\usepackage{parskip}
\usepackage{amsthm}

\usepackage[style=ieee]{biblatex}
\addbibresource{references.bib}

\newcommand{\norm}[1]{\left\lVert#1\right\rVert}
\newcommand{\abs}[1]{\lvert#1\rvert}


\title{Seminar on Graph Theory - Talk 5}
\author{Andres Felipe Jaramillo Garces}
\date{November 2026}

\begin{document}
\maketitle
\section{Origin of the problem}
\cite[\nopp 4.1]{HLM1}
\begin{itemize}
  \item The minimal closed curve problem. which curve encloses the most area given a certai length?
  \item exempify with string
  \item Turn the question around: Given a certain area, what is the shortest closed curve that emcompasses it?
  \item Mention the edge isometric parameter and the vertex isopetrimetric parameter
  \item clarify that the 'area' becomes the number of vertices in a given subset. These two parameters represent two ways to define the perimeter
  \item (write them down) The edge isoperimetic parameter measures the perimeter of an area by counting the outward going edges.
  \item The vertex isoperimetric parameter measures the perimeter by instead counting the nodes adjacent to the subset but external to it. 
  \item Notice that these are different because HLM allows graphs to have multiple edges between them and to themselves, and the adjacency matrix reflects this in its entries.
  \item this problem is in general hard to solve, because although discrete, the obvious answer is just to test all prossible subset that fulfill the requirement.
\end{itemize}
\printbibliography
\section{Construction of the Laplacian}
\begin{itemize}
  \item the incidence matrix: $K_{u,e} = +1$ if $e$ exits $u$, $-1$ if $e$ enters $u$, $0$ otherwise. Its a $V \times E$ matrix.
  \item The gradient: Since a function $f: V \rightarrow R$ on G is a vector, we define $K_{u, e}$ as 1 is the edge $e$ exits $u$ and -1 if $e$ enters u and 0 otherwise. Mind you, this goes back to the interpretation of the edges of the undirected graph as a collection of pairs of directed edges. The gradient is then $fK$ on $E$ (not $V$) as $(fK)_e = f_u - f_v$. 
  \item SHOW EXAMPLE
  \item Let $g: E \rightarrow R$. The divergence maps $g$ to $Kg$, which is indexed by $V$, in particular $(Kg)_v = \sum_{\text{e exits v}} g_e - \sum_{\text{e enters v}} g_e$
  \item SHOW EXAMPLE
  \item The laplacian maps $f: V \rightarrow R$ to $Lf = KK^Tf$. We have that $L_{u,v} = -1$ if $(u,v) \in E$, $deg(u)$ if $u = v$, where $u$ and $v$ are vertices, or 0 otherwise.
  \item you can explain the multiplication $KK^T$ there
  \item $fLf = || fK||^2 = \sum_{(u, v) \in E} (f_u - f_v)^2$
  \item further analysis reveals that $L = dI - A_G$. The paper is ambiguous here with its definitions since it previously defined graphs as allowin multiple edges and self-edges. thi new definition would seem to imply the contrary, but whatever.
\end{itemize}
\section{The central result}
\begin{itemize}
  \item without going into too much detail, there is a result involving the real laplacian and an analogue to the isomeripetric quantities discussed above which relates said analogue with the eigenvalue of the laplacian.
  \item $\frac{d - \lambda}{2} \leq h(G) \leq \sqrt{2d(d -\lambda)}$, there are examples proving both bounds are tight
\end{itemize}

\subsection{Large spectral gap implies high expansion}
We want to prove the lower bound. First rearrage $\lambda \geq d - 2h(G)$. Before processing we recall that the first eigenvalue of the adjacency matrix is $1/\sqrt{n}$, and the first eigenvector is the all-ones vector \textbf{1}. Instead of setting a bound with $\lambda$, we want to prove this using a different relation. First we pick $f \perp 1$. The reason why will become clear soon. We want to find an $f$ such that the Rayleigh Coefficient fulfills $fAf^T/\norm{f}^2 \geq d - 2h(G)$. Why? 

\begin{align*}
  f &= \sum^n_{i=1} c_i v_i = \sum^n_{i=2} c_i v_i \\
  \intertext{Since $f$ was chosen as being perpendicular to \textbf{1}} \\
  \frac{fAf^T}{\norm{f}^2} &= \frac{\sum^n_{i=2} c_i^2 \lambda_i}{\norm{f}^2} = \frac{\sum^n_{i=2} c_i^2 \lambda_i}{\sum^n_{i=2} c_i^2} \leq \frac{\lambda\sum^n_{i=2} c_i^2}{\sum^n_{i=2} c_i^2} = \lambda
\end{align*}

The vector we pick is $f = \abs{\overline{S}}1_S - \abs{S}1_{\overline{S}}$, where $1_S$ is the characteristic vector of a subset $S$ which fulfills $h(G) = \abs{E({S, \overline{S})}}/\abs{S}$ with $\abs{S} \leq n/2$.

Notice that $1_S\cdot1_{\overline{S}} = \sum_{v \in V} (1_X)(1_Y)_v = \sum_{x \in X \cap Y} 1 = \abs{X \cap Y}$

And $1_SA1_S^t = \sum_{u \in S}\sum_{v \in S}A_{uv} = 2\abs{E(S)}$.

And $1_SA1_{\overline{S}}^t = \sum_{u \in S}\sum_{v \in \overline{S}}A_{uv} = 2\abs{E(S, \overline{S})}$.

\subsection{High expansion implies large spectral gap }
To prove the upper bound, we need only find an $S$ for which it holds that $\frac{E(S, \overline{S})}{\abs{S}}$ is smaller than we want it. Thus we want to minimize the numerator to improve our chances.

We pick the eigenvector $g$ to the 2nd eigenvalue $\lambda$. We note that $g \perp v_1 =$ \textbf{1}. We obtain


\begin{align*}
  fLf^t = f(dI - A)f^t = d(ff^t) - fAf^t
\end{align*}

Let $G = (V, E)$ be a finite, connected, $d$-regular undirected graph with $\vert{}V\vert{} = n$ and edge expansion ratio $h = h(G) = \min_{S \subseteq V,\, 0 < \vert{}S\vert{} \le n/2} \frac{\vert{}E(S, \overline{S})\vert{}}{\vert{}S\vert{}}$. From Section 4.5.2, the vector $f = g^+ \in \mathbb{R}^V$ satisfies $f_v = \max(g_v, 0) \ge 0$ for all $v \in V$, and its support $V^+ = \{v \in V : f_v > 0\}$ satisfies $\vert{}V^+\vert{} \le n/2$.

The vertices of $G$ are indexed as $V = \{1, 2, \dots, n\}$ in non-increasing order of the coordinates of $f$:


$$f_1 \ge f_2 \ge \dots \ge f_n \ge 0$$

. Letting $k = \vert{}V^+\vert{}$, this ordering implies:

* $V^+ = \{1, 2, \dots, k\}$ with $f_1 \ge f_2 \ge \dots \ge f_k > 0$,
* $V \setminus V^+ = \{k+1, \dots, n\}$ with $f_{k+1} = f_{k+2} = \dots = f_n = 0$,
* $k \le n/2 \le n - 1$ (for $n \ge 2$), so $V^+ \subseteq \{1, \dots, n-1\}$ and $f_{k+1} = 0$ is well-defined.



For each $i \in \{1, \dots, n-1\}$, the level set is defined as $[i] = \{1, \dots, i\}$, with complement $\overline{[i]} = V \setminus [i] = \{i+1, \dots, n\}$. The quantity $B_f$ is defined by


$$B_f = \sum_{(x,y) \in E} \vert{}f_x^2 - f_y^2\vert{}$$

.

---

The derivation establishes the chain


$$B_f \stackrel{(1)}{=} \sum_{(x,y) \in E,\, x < y} (f_x^2 - f_y^2) \stackrel{(2)}{=} \sum_{(x,y) \in E,\, x < y} \sum_{i=x}^{y-1} (f_i^2 - f_{i+1}^2) \stackrel{(3)}{=} \sum_{i=1}^{n-1} (f_i^2 - f_{i+1}^2) \cdot \vert{}E([i], \overline{[i]})\vert{}$$

$$\stackrel{(4)}{=} \sum_{i \in V^+} (f_i^2 - f_{i+1}^2) \cdot \vert{}E([i], \overline{[i]})\vert{} \stackrel{(5)}{\ge} h \sum_{i \in V^+} (f_i^2 - f_{i+1}^2) \cdot i \stackrel{(6)}{=} h \sum_{i \in V^+} f_i^2 \stackrel{(7)}{=} h \cdot \Vert{}f\Vert{}^2$$

.

**1. Removal of the Absolute Value**
Each undirected non-loop edge $\{u, v\} \in E$ has a unique ordered representation $(x, y)$ with $x < y$ (and any self-loop with $x = y$ contributes $\vert{}f_x^2 - f_x^2\vert{} = 0$). Because $x < y$ implies $f_x \ge f_y \ge 0$, and $t \mapsto t^2$ is monotonically non-decreasing on $[0, \infty)$, it follows that $f_x^2 - f_y^2 \ge 0$ and hence $\vert{}f_x^2 - f_y^2\vert{} = f_x^2 - f_y^2$.

**2. Telescoping Expansion**
For any integers $1 \le x < y \le n$, writing $f_x^2 - f_y^2$ as a telescoping sum of consecutive differences gives:


$$\sum_{i=x}^{y-1} (f_i^2 - f_{i+1}^2) = (f_x^2 - f_{x+1}^2) + (f_{x+1}^2 - f_{x+2}^2) + \dots + (f_{y-1}^2 - f_y^2) = f_x^2 - f_y^2$$

.

**3. Interchange of Summation (Double Counting by Level-Set Cuts)**
For any $1 \le x < y \le n$ and $1 \le i \le n-1$, the index condition $x \le i \le y - 1$ is equivalent to $x \le i$ and $y \ge i + 1$, which is precisely the condition $x \in [i]$ and $y \in \overline{[i]}$ (i.e., the edge $(x, y)$ belongs to the cut $E([i], \overline{[i]})$). Introducing the indicator $\mathbf{1}_{\{x \le i \le y-1\}}$ and interchanging the order of summation yields:


$$\begin{aligned} \sum_{(x,y) \in E,\, x < y} \sum_{i=x}^{y-1} (f_i^2 - f_{i+1}^2) &= \sum_{(x,y) \in E,\, x < y} \sum_{i=1}^{n-1} (f_i^2 - f_{i+1}^2) \cdot \mathbf{1}_{\{x \le i \le y-1\}} \\ &= \sum_{i=1}^{n-1} (f_i^2 - f_{i+1}^2) \sum_{(x,y) \in E,\, x < y} \mathbf{1}_{\{x \in [i],\, y \in \overline{[i]}\}} \\ &= \sum_{i=1}^{n-1} (f_i^2 - f_{i+1}^2) \cdot \vert{}E([i], \overline{[i]})\vert{} \end{aligned}$$

.

**4. Restriction of the Summation Domain to $V^+$**
Since $V^+ = \{1, \dots, k\}$ and $f_j = 0$ for all $j \in \{k+1, \dots, n\}$, every index $i \in \{k+1, \dots, n-1\}$ satisfies both $f_i = 0$ and $f_{i+1} = 0$, so $f_i^2 - f_{i+1}^2 = 0$. Dropping these vanishing terms restricts the summation from $\{1, \dots, n-1\}$ to $V^+ = \{1, \dots, k\}$.

**5. Application of the Edge Expansion Ratio**
For every $i \in V^+ = \{1, \dots, k\}$, the level set $[i]$ has cardinality $\vert{}[i]\vert{} = i$ satisfying $0 < \vert{}[i]\vert{} \le \vert{}V^+\vert{} \le n/2$. By the definition of $h = h(G)$:


$$\frac{\vert{}E([i], \overline{[i]})\vert{}}{\vert{}[i]\vert{}} \ge h \implies \vert{}E([i], \overline{[i]})\vert{} \ge h \cdot i.$$


Because $f_i \ge f_{i+1} \ge 0$, each factor $(f_i^2 - f_{i+1}^2)$ is non-negative. Multiplying the inequality $\vert{}E([i], \overline{[i]})\vert{} \ge h \cdot i$ by $(f_i^2 - f_{i+1}^2) \ge 0$ and summing over all $i \in V^+$ preserves the inequality:


$$\sum_{i \in V^+} (f_i^2 - f_{i+1}^2) \cdot \vert{}E([i], \overline{[i]})\vert{} \ge \sum_{i \in V^+} (f_i^2 - f_{i+1}^2) \cdot h i = h \sum_{i \in V^+} (f_i^2 - f_{i+1}^2) \cdot i$$

.

**6. Summation by Parts (Collapsing the Telescopic Sum)**
Reindexing the second term of $\sum_{i=1}^{k} i (f_i^2 - f_{i+1}^2)$ via $j = i + 1$:


$$\begin{aligned} \sum_{i=1}^{k} i (f_i^2 - f_{i+1}^2) &= \sum_{i=1}^{k} i f_i^2 - \sum_{i=1}^{k} i f_{i+1}^2 \\ &= \sum_{i=1}^{k} i f_i^2 - \sum_{j=2}^{k+1} (j - 1) f_j^2 \\ &= 1 \cdot f_1^2 + \sum_{i=2}^{k} \bigl(i - (i - 1)\bigr) f_i^2 - k f_{k+1}^2 \\ &= \sum_{i=1}^{k} f_i^2 - k f_{k+1}^2. \end{aligned}$$


Because $k = \vert{}V^+\vert{}$ implies $k + 1 \notin V^+$, we have $f_{k+1} = 0$, so the boundary term $-k f_{k+1}^2$ vanishes, leaving $h \sum_{i \in V^+} f_i^2$.

**7. Identification with $\Vert{}f\Vert{}^2$**
Since $f_i = 0$ for all $i \in V \setminus V^+$,


$$\Vert{}f\Vert{}^2 = \sum_{i \in V} f_i^2 = \sum_{i \in V^+} f_i^2 + \sum_{i \in V \setminus V^+} f_i^2 = \sum_{i \in V^+} f_i^2,$$


yielding $h \sum_{i \in V^+} f_i^2 = h \cdot \Vert{}f\Vert{}^2$.


\end{document}

